\section{Choosing among different optimization strategies }
\label{Decision tree}
There are several considerations which determine whether full-AD, HG or semi-AD is the optimal strategy for a given optimization problem, see the decision tree Fig.\ \ref{fig:decision tree}. Whenever the Hilbert space dimension $d$ and the number of time steps $N$ are sufficiently small, then memory usage may not be a bottleneck. In this case, full-AD may be preferable, since it eliminates the need for hard-coded gradients. If the memory cost of full-AD is not affordable but analytical gradients are tedious to compute, semi-AD may be a viable compromise, though still less memory efficient than HG. In all other cases, HG is the method of choice, since it can handle optimization with larger $N$ and $d$ compared to full-AD and semi-AD. 

The choice of implementing HG via scaling and squaring or matrix diagonalization should be informed by whether the Hamiltonian is sparse and $\kappa$ scales better than $\Theta(d^2)$. In the latter case, scaling and squaring is preferable and leads to an overall memory usage $\sim \Theta(d+\kappa)$, which is better than the memory scaling $O(d^2)$ of matrix diagonalization. By contrast, for Hamiltonians with $\kappa\sim \Theta(d^2)$, these two methods are equivalent in memory usage scaling and we are free to select the one with better runtime for a specific optimization task. While the runtime for optimization using matrix diagonalization scales as $O(d^3)$ [\ref{scalingmatrixdiag}], the runtime for optimization based on scaling and squaring differs between state-transfer and gate optimizations: (1) for state transfer, the runtime scales as $\Theta(\mu d^2)$; hence, whenever $\mu$ scales better than $\Theta(d)$, scaling and squaring wins over matrix diagonalization. (2) for gate optimization, the runtime scales as $\Theta(\mu d^3)$; hence matrix diagonalization is preferable. 


